\chapter{Localization Analysis}
\label{appendix:localization}

Both localizers of \autoref{sec:priming-localizers} use the primitive $\textsc{Measure}(t)$, which
runs the batch $\batch{\spl{t}}$ and returns the probe reading $r(t)$. We show each localizer is
optimal for what it returns and summarize the trade-off in \autoref{tab:loc-compare}.

\section{Linear scan}
\label{app:loc-scan}

The scan (\autoref{alg:prime-scan}) walks the split backwards from the probe end. From the known
$r(k+1)$ it measures $r(k),r(k-1),\dots$ and stops at the first neighbor that disagrees. That split
is a boundary, and its column is the highest leaking column.

\begin{proposition}[Cost of the scan]
\label{prop:scancost}
The linear scan returns $\tmax=\max\{t:r(t)\neq r(t+1)\}$ and runs at most $k$ batches.
\end{proposition}
\begin{proof}
It stops at the first $t$, counting down from $k$, with $r(t)\neq r(t+1)$. Since it did not stop
earlier, $r(t+1)=\dots=r(k+1)$, so no split above $t$ is a boundary and $t=\tmax$. The interior
splices it runs lie among $\spl{k},\dots,\spl{1}$, at most $k$.
\end{proof}

\section{Binary search}
\label{app:loc-bsearch}

Binary search (\autoref{alg:prime-bsearch}) keeps a window $[\ell,h]$ whose end readings differ, so
a boundary lies inside, and halves it until it spans one split. \autoref{fig:binsearch} runs it on
$12$ inputs with the leak at column~$4$.

\begin{proposition}[Binary search returns a boundary]
\label{prop:bsearch}
Binary search makes at most $\lceil\log_2(k+1)\rceil$ runs and returns a boundary. By the argument
of \autoref{sec:priming-boundary} its pair is a violation.
\end{proposition}
\begin{proof}
The invariant $r(\ell)\neq r(h)$ holds initially and is preserved: if $r(m)\neq r(\ell)$ then
$[\ell,m]$ has unequal ends, otherwise $r(m)=r(\ell)\neq r(h)$ and $[m,h]$ does. Each step halves
$h-\ell$, so from $k+1$ it reaches $1$ in at most $\lceil\log_2(k+1)\rceil$ runs, ending with
$h=\ell+1$ and $r(\ell)\neq r(\ell+1)$.
\end{proof}

\begin{figure}[p]
\centering
\fitwidth{\newcommand{\bsrow}[4]{%
  \node[font=\scriptsize,anchor=east] at (-0.45,#1) {#3};
  \node[ci] at (0,#1) {IS};
  \foreach \c in {0,...,10}{
    \pgfmathsetmacro{\x}{0.68*(\c+1)}
    \ifnum\c<#2 \node[ca] at (\x,#1) {$\inp{\c}{a}$};
    \else       \node[cb] at (\x,#1) {$\inp{\c}{b}$}; \fi}
  \ifnum#2=12 \node[pa] at (8.16,#1) {$\inp{11}{a}$};
  \else       \node[pb] at (8.16,#1) {$\inp{11}{b}$}; \fi
  \node[font=\scriptsize,anchor=west] at (8.6,#1) {#4};}
\newcommand{\bsneq}[2]{%
  \draw[culprit,line width=1pt] (10.05,#1) -- (10.25,#1) -- (10.25,#2) -- (10.05,#2);
  \node[culprit,font=\small,anchor=west] at (10.25,{(#1+#2)/2}) {$\neq$};}
\newcommand{\bswin}[5]{%
  \draw[culprit,line width=1.1pt] ({0.68*(#3+1)-0.37},#1+0.35) rectangle ({0.68*#4+0.37},#2-0.35);
  \node[culprit,font=\scriptsize,anchor=south] at ({0.34*(#3+#4)+0.34},#1+0.35) {#5};}
\begin{tikzpicture}[font=\small,
  ce/.style={draw,minimum width=0.62cm,minimum height=0.52cm,font=\tiny,inner sep=0pt},
  ca/.style={ce,fill=softA}, cb/.style={ce,fill=softB},
  ci/.style={ce,fill=iscol!25,font=\tiny\bfseries},
  pa/.style={ce,line width=1pt,draw=probe,fill=softA},
  pb/.style={ce,line width=1pt,draw=probe,fill=softB}]
  \node[font=\scriptsize,anchor=east] at (-0.45,0.75) {column};
  \foreach \c in {0,...,10}{\node[font=\scriptsize] at ({0.68*(\c+1)},0.75) {\c};}
  \node[probe,font=\scriptsize] at (8.16,0.75) {11};
  \node[probe,font=\scriptsize,anchor=west] at (8.6,0.75) {(probe)};
  \begin{scope}[yshift=0cm]
  \node[font=\footnotesize\bfseries,anchor=west] at (-3.6,0.1) {Iteration 1};
  \bsrow{-0.4}{12}{$h{=}12$: $\batch{\spl{12}}$ (known)}{\textcolor{laneA}{$r(12)=\rho_a$}}
  \bsrow{-1.1}{6}{$m{=}6$: $\batch{\spl{6}}$ (run)}{\textcolor{laneA}{$r(6)=\rho_a$}}
  \bsrow{-1.8}{0}{$\ell{=}0$: $\batch{\spl{0}}$ (known)}{\textcolor{laneB}{$r(0)=\rho_b$}}
  \bswin{-0.4}{-1.8}{0}{12}{window $[0,12]$: columns $0..11$}
  \bsneq{-1.1}{-1.8}
  \node[font=\scriptsize,anchor=west] at (-3.6,-3.0) {$r(6)\neq r(0)$: the change is in columns $0..5$, between $\spl{6}$ and $\spl{0}$ $\Rightarrow$ go left, $h:=6$};
  \end{scope}
  \begin{scope}[yshift=-3.75cm]
  \node[font=\footnotesize\bfseries,anchor=west] at (-3.6,0.1) {Iteration 2};
  \bsrow{-0.4}{6}{$h{=}6$: $\batch{\spl{6}}$ (known)}{\textcolor{laneA}{$r(6)=\rho_a$}}
  \bsrow{-1.1}{3}{$m{=}3$: $\batch{\spl{3}}$ (run)}{\textcolor{laneB}{$r(3)=\rho_b$}}
  \bsrow{-1.8}{0}{$\ell{=}0$: $\batch{\spl{0}}$ (known)}{\textcolor{laneB}{$r(0)=\rho_b$}}
  \bswin{-0.4}{-1.8}{0}{6}{window $[0,6]$: columns $0..5$}
  \bsneq{-0.4}{-1.1}
  \draw[decorate,decoration={brace,amplitude=3pt,mirror},laneB] ({0.68*(6+1)-0.33},-2.2) -- ({0.68*(11+1)+0.33},-2.2) node[midway,below=1pt,font=\scriptsize] {suffix: lane $b$ in all three};
  \node[font=\scriptsize,anchor=west] at (-3.6,-3.0) {$r(3)=r(0)$, so $r(3)\neq r(6)$: the change is in columns $3..5$, between $\spl{6}$ and $\spl{3}$ $\Rightarrow$ go right, $\ell:=3$};
  \end{scope}
  \begin{scope}[yshift=-7.5cm]
  \node[font=\footnotesize\bfseries,anchor=west] at (-3.6,0.1) {Iteration 3};
  \bsrow{-0.4}{6}{$h{=}6$: $\batch{\spl{6}}$ (known)}{\textcolor{laneA}{$r(6)=\rho_a$}}
  \bsrow{-1.1}{4}{$m{=}4$: $\batch{\spl{4}}$ (run)}{\textcolor{laneB}{$r(4)=\rho_b$}}
  \bsrow{-1.8}{3}{$\ell{=}3$: $\batch{\spl{3}}$ (known)}{\textcolor{laneB}{$r(3)=\rho_b$}}
  \bswin{-0.4}{-1.8}{3}{6}{window $[3,6]$: columns $3..5$}
  \bsneq{-0.4}{-1.1}
  \draw[decorate,decoration={brace,amplitude=3pt,mirror},laneA] ({0.68*(0+1)-0.33},-2.2) -- ({0.68*(2+1)+0.33},-2.2) node[midway,below=1pt,font=\scriptsize] {prefix: lane $a$};
  \draw[decorate,decoration={brace,amplitude=3pt,mirror},laneB] ({0.68*(6+1)-0.33},-2.2) -- ({0.68*(11+1)+0.33},-2.2) node[midway,below=1pt,font=\scriptsize] {suffix: lane $b$};
  \node[font=\scriptsize,anchor=west] at (-3.6,-3.0) {$r(4)=r(3)$, so $r(4)\neq r(6)$: the change is in columns $4,5$, between $\spl{6}$ and $\spl{4}$ $\Rightarrow$ go right, $\ell:=4$};
  \end{scope}
  \begin{scope}[yshift=-11.25cm]
  \node[font=\footnotesize\bfseries,anchor=west] at (-3.6,0.1) {Iteration 4};
  \bsrow{-0.4}{6}{$h{=}6$: $\batch{\spl{6}}$ (known)}{\textcolor{laneA}{$r(6)=\rho_a$}}
  \bsrow{-1.1}{5}{$m{=}5$: $\batch{\spl{5}}$ (run)}{\textcolor{laneA}{$r(5)=\rho_a$}}
  \bsrow{-1.8}{4}{$\ell{=}4$: $\batch{\spl{4}}$ (known)}{\textcolor{laneB}{$r(4)=\rho_b$}}
  \bswin{-0.4}{-1.8}{4}{6}{window $[4,6]$: columns $4,5$}
  \bsneq{-1.1}{-1.8}
  \draw[decorate,decoration={brace,amplitude=3pt,mirror},laneA] ({0.68*(0+1)-0.33},-2.2) -- ({0.68*(3+1)+0.33},-2.2) node[midway,below=1pt,font=\scriptsize] {prefix: lane $a$};
  \draw[decorate,decoration={brace,amplitude=3pt,mirror},laneB] ({0.68*(6+1)-0.33},-2.2) -- ({0.68*(11+1)+0.33},-2.2) node[midway,below=1pt,font=\scriptsize] {suffix: lane $b$};
  \node[font=\scriptsize,anchor=west] at (-3.6,-3.0) {$r(5)\neq r(4)$: the change is in column $4$, between $\spl{5}$ and $\spl{4}$ $\Rightarrow$ go left, $h:=5$};
  \end{scope}
  \begin{scope}[yshift=-15.0cm]
  \node[font=\footnotesize\bfseries,anchor=west] at (-3.6,0.1) {Stop: $h=\ell+1$};
  \bsrow{-0.4}{5}{$h{=}5$: $\batch{\spl{5}}$ (known)}{\textcolor{laneA}{$r(5)=\rho_a$}}
  \bsrow{-1.1}{4}{$\ell{=}4$: $\batch{\spl{4}}$ (known)}{\textcolor{laneB}{$r(4)=\rho_b$}}
  \draw[culprit,line width=1.1pt] ({0.68*5-0.33},-0.4+0.31) rectangle ({0.68*5+0.33},-1.1-0.31);
  \node[culprit,font=\scriptsize,anchor=south] at ({0.68*5},-0.4+0.31) {window $[4,5]$: column $4$ only};
  \bsneq{-0.4}{-1.1}
  \draw[decorate,decoration={brace,amplitude=3pt,mirror},laneA] ({0.68*(0+1)-0.33},-1.5) -- ({0.68*(3+1)+0.33},-1.5) node[midway,below=1pt,font=\scriptsize] {prefix: lane $a$};
  \draw[decorate,decoration={brace,amplitude=3pt,mirror},laneB] ({0.68*(5+1)-0.33},-1.5) -- ({0.68*(11+1)+0.33},-1.5) node[midway,below=1pt,font=\scriptsize] {suffix: lane $b$};
  \node[font=\scriptsize,anchor=west] at (-3.6,-2.35) {the two rows differ only in column $4$, and $r(5)\neq r(4)$: boundary $t=4$, column $4$ leaks, four runs in total.};
  \end{scope}
\end{tikzpicture}
}
\caption{Binary search on $12$ inputs with the leak at column~$4$. The red frame marks the window
$[\ell,h]$, columns $\ell,\dots,h-1$.}
\label{fig:binsearch}
\end{figure}

\begin{table}[t]
\centering\small
\setlength{\tabcolsep}{8pt}
\caption{The two localizers, per direction and per repetition. There are $k+1$ columns, and $\tmax$
is the highest leaking column.}
\label{tab:loc-compare}
\begin{tabular}{@{}lll@{}}
\toprule
 & Linear scan & Binary search \\
\midrule
Runs (worst case) & $k$ & $\lceil\log_2(k{+}1)\rceil$ \\
Lower bound & $k$~(\autoref{thm:lb-scan}) & $\lceil\log_2(k{+}1)\rceil$~(\autoref{thm:lb-bsearch}) \\
Returns & $\tmax$ (all if continued) & one boundary \\
\bottomrule
\end{tabular}
\end{table}

\section{Lower bounds}
\label{app:loc-lower}

Both localizers are optimal for what they return. The arguments use the probe $\textsc{Measure}(t)$
and a family of \emph{step sequences}. Fix two distinct readings $\rho_a\neq\rho_b$. For each
$j\in\{0,\dots,k\}$, the step sequence $S_j$ is the reading pattern a single leaking column $j$
produces:
\[
  \textsc{Measure}(t)=\begin{cases}\rho_b & t\le j,\\ \rho_a & t>j.\end{cases}
\]
Its only boundary is $j$, so $\tmax=j$. A localizer is deterministic, so the same readings yield the
same runs and the same output.

\begin{theorem}[Finding $\tmax$ needs $k$ runs]
\label{thm:lb-scan}
Every localizer that always returns $\tmax$ runs at least $k$ batches on some reading sequence, so
the linear scan is optimal.
\end{theorem}
\begin{proof}
Suppose the hardware behaves as $S_j$ and take any localizer that always returns $\tmax$. Suppose it
never runs some split $u$ with $j<u\le k$. Build $S'$ identical to $S_j$ except that
$\textsc{Measure}(u)=\gamma$ for a third value $\gamma\neq\rho_a,\rho_b$. On every split the
localizer runs, $S'$ and $S_j$ agree, so it makes the identical runs and emits the same output $j$
on $S'$. But $S'$ has a boundary at $u>j$, so $\tmax\ge u>j$ and the output is wrong. Hence the
localizer runs every split in $\{j{+}1,\dots,k\}$. Taking $j=0$ forces all $k$ splits, matching the
linear scan.
\end{proof}

\begin{theorem}[Finding a boundary needs $\lceil\log_2(k+1)\rceil$ runs]
\label{thm:lb-bsearch}
Every localizer that always returns a boundary runs at least $\lceil\log_2(k+1)\rceil$ batches on
some step sequence, so binary search is optimal.
\end{theorem}
\begin{proof}
The $k+1$ step sequences $S_0,\dots,S_k$ are the candidates, and the localizer must output the
correct $j$. A probe at split $t$ returns $\rho_b$ when $j\ge t$ and $\rho_a$ otherwise, so it
reveals only whether $j\ge t$. An adversary answers the probes without fixing $j$ in advance,
keeping its answers consistent across probes by retaining, after each probe, the larger of the two
candidate sets it leaves open. Each probe then removes at most half the candidates. The localizer
cannot stop while two candidates remain, since both are consistent with every answer so far yet
have different boundaries. Driving $k+1$ candidates down to one therefore takes at least
$\lceil\log_2(k+1)\rceil$ probes, which binary search attains.
\end{proof}
